Le projet ITInéraire a permis de concevoir et d'implémenter un système de localisation embarqué robuste, capable de compenser les faiblesses d'un positionnement GNSS standard. En fusionnant les données satellites avec celles de l'IMU, nous avons pu générer une estimation de trajectoire continue et physiquement cohérente. 

Ce projet a été particulièrement intéressant d’un point de vue conception et intégration matérielle. Se rendre compte du fonctionnement réel et des limites des capteurs physiques (dérive de l'IMU, sauts du GNSS) est la clé pour comprendre comment développer et paramétrer un algorithme de fusion de données robuste comme le Filtre de Kalman Étendu.

De plus, ce projet a été l’occasion de confronter nos différents points de vue sur certains aspects conceptuels et mathématiques, notamment lors de l'abandon du magnétomètre ou du réglage des matrices de covariance. Mettre en lumière nos différentes façons de penser, se projeter sur les fonctionnalités interactives à implémenter pour la résolution du GeoGuessr, et la communication au sein du groupe ont permis de surmonter les obstacles techniques et de concevoir ce projet de manière réfléchie et réussie.